\documentclass[10pt,a4paper]{amsart}
\usepackage[margin=19mm]{geometry}
\usepackage{amsmath,amssymb,mathtools}
\usepackage[scr=rsfs,cal=euler]{mathalfa}
\usepackage{xfrac}
\usepackage[unicode,hidelinks,bookmarks=false]{hyperref}
\newcommand{\ie}{i.e.\ }
\newcommand{\cf}{cf.\ }
\newcommand{\resp}{respectively\ }
\newcommand{\Subsection}{\subsection}
\providecommand{\nobreakdash}{\nobreak\hbox{-}}
\setlength{\emergencystretch}{2em}
\multlinegap0pt
\usepackage{bm}
\usepackage{bbm}
\usepackage{dsfont}
\usepackage{xspace}
\usepackage{cancel}
\usepackage{mathtools}
\usepackage{amsthm}
\makeatletter
\@ifundefined{theorem}{\newtheorem{theorem}{Theorem}}{}
\@ifundefined{proposition}{\newtheorem{proposition}[theorem]{Proposition}}{}
\makeatother
\newcommand{\mcI}{\mathcal{I}}
\title[Circuit complexity lower bounds for quantum spin glasses]{Circuit complexity lower bounds for quantum spin glasses}
\author{Omar Al-Ghattas, David Gamarnik}
\date{Source: arXiv:2607.14384v2 (CC BY 4.0)}
\begin{document}
\maketitle

\noindent\emph{Source and licence.} Circuit complexity lower bounds for quantum spin glasses. Version arXiv:2607.14384v2 (CC BY 4.0), distributed under the Creative Commons Attribution 4.0 International licence, as recorded in the arXiv licence metadata for this exact version and shown on the abstract page https://arxiv.org/abs/2607.14384v2. Full rights notice: licenses/arxiv-2607.14384-CC-BY-4.0.txt.\par\smallskip

\noindent\textit{Editorial scope.} Counted unit: the Proposition whose source label is \texttt{prop:full-residual-entropy} in the article (source0.tex lines 1897--1910, source label \texttt{prop:full-residual-entropy}), delivered together with the article's own complete original proof of it (source0.tex lines 1912--2278, one \texttt{proof} environment of 367 lines). No mathematics is written, repaired or re-derived: statement and proof are the article's own text line for line. Required context retained uncounted: none. Every definition, display and label of those lines is retained unchanged. Omitted from this package and not redistributed: 1--1896, 1911--1911, 2279--5813. Nothing inside a retained excerpt is omitted or altered: every retained body is delivered exactly as the article's lines are written, so no omission receipt of any kind accompanies this package. Citations: the retained excerpts carry no citation command at all, so no bibliography entry of the article is transcribed and no reference is redistributed. Reference adaptation: none was needed and none was made; every reference inside the delivered unit resolves inside it, so no unresolved reference or citation is produced. Printed slips kept literal: no printed word, symbol, hypothesis or number of the counted statement or of its proof was altered. Layout and class: the article's own class is \texttt{article} with packages including subfigure, inputenc, mmap, fontenc, authblk, latexsym, graphics, amsmath, xspace, amssymb, psfrag, pst-all, color, bbm; the wrapper is the lane's pinned \texttt{amsart} editorial wrapper at 10pt on A4, which restates the theorem-like environments the retained text needs and adds the article's own macro definitions that the retained text needs (\texttt{\textbackslash mcI}), with extra wrapper packages (bm, bbm, dsfont, xspace, cancel, mathtools). The delivered numbering is the unit's own. Assets: no retained span contains a figure environment, a graphics-inclusion command, a PDF-page inclusion, a table or a TikZ picture; no image file is copied into this package, the package declares no asset dependency and no image file is redistributed with it. Licence: version arXiv:2607.14384v2 of this article is distributed under the Creative Commons Attribution 4.0 International licence, as recorded in the arXiv licence metadata for this exact version and shown on the abstract page https://arxiv.org/abs/2607.14384v2. A full rights notice is delivered at licenses/arxiv-2607.14384-CC-BY-4.0.txt. Characters: every retained excerpt is pure ASCII, so the delivered engine, XeLaTeX with its default Latin Modern text font, reports no missing character.
\begin{proposition}[Entropy bound, mean-field regime]
\label{prop:full-residual-entropy}
Fix \(p\ge2\). Let \(r=r_n\ge0\). Let \(\mathsf{d}_{R}\) be the canonical metric of the residual process
$R$ on \(\mathsf{T}_{n,r}^{(D)}\). Then, for every \(\eta>0\),
\[
\log N(\mathsf{T}_{n,r}^{(D)},\mathsf{d}_{R},\eta)
\le
C_pDn2^D(\log n+D)
+
C_pDn2^D
\log_+\left(\frac{C_pDn2^D}{\eta}\right),
\]
uniformly over \(r\ge0\).
\end{proposition}

\begin{proof}
We identify a parameter
\(\theta\in\mathsf{T}_{n,r}^{(D)}\) with its corresponding circuit unitary \(U\) when no
confusion can arise. Thus $U=V_DV_{D-1}\cdots V_1$, where each layer has the form
\[
    V_t
    =
    \bigotimes_{e\in M_t}U_e^{(t)},
    \qquad
    U_e^{(t)}\in U(4),
\]
for a perfect matching \(M_t\) on \([n+r]\). Throughout, whenever we
write a tensor product over a partial matching, the resulting operator is
understood to act as the identity on all qubits not covered by that matching.

We prove the entropy bound by replacing each circuit, in the metric
\(\mathsf{d}_{R}\), by an equivalent pruned and relabeled circuit acting on at most \(n2^D\) qubits.

We now construct the pruned circuit. It is enough to track how physical
observables pull back through the circuit. Here and below, a physical observable means an observable supported on the first \(n\) qubits.

Write $U=V_DV_{D-1}\cdots V_1$, where \(V_1\) acts first and \(V_D\) acts last. Given a physical observable \(X\), define the pulled-back observables
\[
    X_D:=X,
    \qquad
    X_{t-1}:=V_t^*X_tV_t,
    \qquad t=D,D-1,\dots,1.
\]
Thus \(X_t\) is the observable obtained after pulling \(X\) backward through
the later layers \(D,D-1,\dots,t+1\). In particular, $X_0=U^*XU.$
Since \(X_D=X\) is physical, \(\operatorname{supp}(X_D)\subseteq[n]\). We will
use a joint support bound that works for all physical observables at once, and therefore initialize $S_D:=[n].$

Now define the backward active sets recursively. For \(t=D,D-1,\dots,1\), set
\[
    R_t:=\{e\in M_t:e\cap S_t\neq\emptyset\},
    \qquad
    S_{t-1}:=S_t\cup\bigcup_{e\in R_t}e.
\]
Here \(S_t\) is a set which contains the support of $X_t$, uniformly over all physical observables \(X\) considered in the residual process. The
set \(R_t\) consists of precisely those gates in layer \(t\) that touch this
current backward support. These are the only gates in layer \(t\) that can
affect the pulled-back observable when passing from time \(t\) to time
\(t-1\).

Since \(M_t\) is a matching, each qubit in \(S_t\) belongs to at most one edge
of \(M_t\). Therefore each active qubit can introduce at most one new qubit
when we pull backward through layer \(t\). Hence $|S_{t-1}|\le 2|S_t|,$ and iterating gives
\[
    |S_t|\le n2^{D-t},
    \qquad t=0,1,\dots,D.
\]
In particular, \(|S_0|\le n2^D\). Moreover, since every edge of \(R_t\)
touches \(S_t\), we have \(|R_t|\le |S_t|\), and hence
\[
    \sum_{t=1}^D |R_t|
    \le
    \sum_{t=1}^D |S_t|
    \le
    n(2^D-1).
\]

Define the pruned layer and pruned circuit by
\[
    V_t^{\mathrm{pr}}
    :=
    \bigotimes_{e\in R_t}U_e^{(t)},
    \qquad 
    U^{\mathrm{pr}}
    :=
    V_D^{\mathrm{pr}}\cdots V_1^{\mathrm{pr}},
\]
where \(V_t^{\mathrm{pr}}\) acts as the identity on all qubits not covered by
\(R_t\). We claim that, for every physical observable \(X\), $U^*XU=(U^{\mathrm{pr}})^*XU^{\mathrm{pr}}$.

To prove the claim, use the pulled-back observables \(X_t\) defined above. We prove by induction that \(\operatorname{supp}(X_t)\subseteq S_t\) for every
\(t=D,D-1,\dots,0\), and that
\[
    X_{t-1}=(V_t^{\mathrm{pr}})^*X_tV_t^{\mathrm{pr}}
    \qquad\text{for }t=D,D-1,\dots,1.
\]
The base case holds because \(X_D=X\) is physical, so
\(\operatorname{supp}(X_D)\subseteq[n]=S_D\).

Assume that \(\operatorname{supp}(X_t)\subseteq S_t\). Let
\[
    V_t^{\mathrm{bad}}
    :=
    \bigotimes_{e\in M_t\setminus R_t}U_e^{(t)}.
\]
Since \(M_t\) is a matching, the gates in \(R_t\) and \(M_t\setminus R_t\)
act on disjoint pairs, and hence $V_t=V_t^{\mathrm{bad}}V_t^{\mathrm{pr}}$. By definition of \(R_t\), every edge in \(M_t\setminus R_t\) is disjoint from
\(S_t\). Hence \(V_t^{\mathrm{bad}}\) is supported outside
\(\operatorname{supp}(X_t)\), and therefore $(V_t^{\mathrm{bad}})^*X_tV_t^{\mathrm{bad}}=X_t$. It follows that
\[
    X_{t-1}
    =
    V_t^*X_tV_t
    =
    (V_t^{\mathrm{pr}})^*
    (V_t^{\mathrm{bad}})^*
    X_t
    V_t^{\mathrm{bad}}
    V_t^{\mathrm{pr}}
    =
    (V_t^{\mathrm{pr}})^*X_tV_t^{\mathrm{pr}} .
\]
Finally, conjugation by \(V_t^{\mathrm{pr}}\) can only spread the support along
the edges in \(R_t\). Thus
\[
    \operatorname{supp}(X_{t-1})
    \subseteq
    S_t\cup\bigcup_{e\in R_t}e
    =
    S_{t-1}.
\]
This proves the induction. Since \(X_0=U^*XU\), the displayed identity for
\(X_{t-1}\) also gives $U^*XU=(U^{\mathrm{pr}})^*XU^{\mathrm{pr}}$, as claimed.

Let \(\theta^{\mathrm{pr}}\) denote the circuit parameter obtained from
\(\theta\) by replacing \(U\) with \(U^{\mathrm{pr}}\), equivalently by setting
all deleted gates to the identity. Applying this identity to the physical
Pauli products \(P_I^a\otimes I_{\mathrm{anc}}\) and to the one-site physical
observables \(\sigma_i^b\otimes I_{\mathrm{anc}}\), we obtain $\chi_{I,a}(\theta)=\chi_{I,a}(\theta^{\mathrm{pr}})$ for every \(I\in\mcI_p^n\) and \(a\in\{1,2,3\}^p\). Consequently, $\mathsf d_R(\theta,\theta^{\mathrm{pr}})=0$.

After pruning, each circuit parameter \(\theta\) is \(\mathsf d_R\)-equivalent
to the pruned parameter \(\theta^{\mathrm{pr}}\), whose nonidentity gates
involve only the qubits in the backward active set \(S_0\). In particular, at
most \(n(2^D-1)\) ancillas can remain active. However, these active ancillas
may have arbitrary labels among the \(r\) available ancillas. If we counted
these labels directly, we would reintroduce an \(r\)-dependent combinatorial
factor. Since all ancillas are initialized in the same state and the observables
in the residual process act only on the physical qubits, we may permute ancilla
labels, while keeping the physical labels fixed, without changing any physical
expectation.

Concretely, let $A(\theta^{\mathrm{pr}}):=S_0\setminus[n]$ be the set of ancilla labels appearing in the joint backward support of the
physical outputs. Since \(|S_0|\le n2^D\), we have $|A(\theta^{\mathrm{pr}})|\le n(2^D-1)$.

Set
\[
    q_*:=\min\{r,n(2^D-1)\},
    \qquad
    Q:=n+q_*.
\]
Then \(Q\le n2^D\). Let
\[
    K:=\{1,\dots,n\}\cup\{n+1,\dots,n+q_*\}
\]
be the set consisting of the physical qubits and the first \(q_*\) ancillas.

Since \(|A(\theta^{\mathrm{pr}})|\le q_*\), there exists a permutation \(\pi\)
of \([n+r]\) such that \(\pi(i)=i\) for every physical label \(i\in[n]\), and
\[
    \pi(A(\theta^{\mathrm{pr}}))\subseteq\{n+1,\dots,n+q_*\}.
\]
Let \(P_\pi\) denote the corresponding qubit-permutation unitary, and define $\widetilde U:=P_\pi U^{\mathrm{pr}}P_\pi^*$.

Let \(\widetilde\theta\) denote the corresponding relabeled circuit parameter.
Then \(\widetilde\theta\) is a depth-\(D\) matching-circuit parameter whose
nonidentity gates are all supported on \(K\).

We now check that this relabeling preserves all residual coefficients. Since
\(\pi\) fixes the physical labels, for every physical observable \(X\), $P_\pi^*XP_\pi=X$. Also \(P_\pi|0^{n+r}\rangle=|0^{n+r}\rangle\). Therefore
\[
\langle0^{n+r}|\widetilde U^*X\widetilde U|0^{n+r}\rangle
=
\langle0^{n+r}|
U^{\mathrm{pr}*}XU^{\mathrm{pr}}
|0^{n+r}\rangle.
\]
Applying this identity to all physical Pauli products and one-site physical
Pauli observables gives
\[
    \chi_{I,a}(\widetilde\theta)
    =
    \chi_{I,a}(\theta^{\mathrm{pr}})
    =
    \chi_{I,a}(\theta)
\]
for every \(I\in\mcI_p^n\) and \(a\in\{1,2,3\}^p\). Thus $\mathsf d_R(\theta,\widetilde\theta)=0$.


Consequently, every circuit parameter in \(\mathsf{T}_{n,r}^{(D)}\) is
\(\mathsf{d}_{R}\)-equivalent to a circuit parameter whose nonidentity gates act only
on a canonical set of size $Q\le n2^D$. 


It remains to cover depth-\(D\) matching circuits on the set \(K\). Let \(N_Q\) be the number of partial matchings on a \(Q\)-element set. A partial matching with exactly \(m\) edges is
obtained by choosing \(2m\) vertices and partitioning them into \(m\) unordered
pairs. Hence
\[
    N_Q
    =
    \sum_{m=0}^{\lfloor Q/2\rfloor}
    \binom{Q}{2m} \frac{(2m)!}{2^m m!}
    =
    \sum_{m=0}^{\lfloor Q/2\rfloor}
    \frac{Q!}{(Q-2m)!2^m m!}.
\]
For each \(m\),
\[
    \frac{Q!}{(Q-2m)!2^m m!}
    \le Q^{2m}\le Q^Q.
\]
Thus
\[
    N_Q
    \le
    (Q+1)Q^Q
    \le
    \exp (CQ\log(eQ)).
\]
A depth-\(D\) architecture is a sequence of \(D\) partial matchings on \(K\). Therefore the number of possible architectures is at most \(N_Q^D\). Hence
\[
    \log(N_Q^D)
    =
    D\log N_Q
    \le
    C DQ\log(eQ).
\]
Since \(Q\le n2^D\), we have
\[
    DQ\log(eQ)
    \le
    Dn2^D\log(en2^D)
    \le
    C Dn2^D(\log n+D).
\]
Thus the logarithm of the number of possible depth-\(D\) architectures is at most $C Dn2^D(\log n+D)$.

We now cover the remaining unitary degrees of freedom on the active set \(K\). The counting is done in two stages. We first fix the pruned, relabeled architecture, namely the sequence of partial matchings \(\mathcal M=(M_1,\dots,M_D)\) on \(K\). Then we discretize the two-qubit unitaries assigned to the edges of these matchings. The nets obtained for the different architectures will later be combined by taking a union over the possible choices of \(\mathcal M\).

Fix a depth-\(D\) architecture $\mathcal M=(M_1,\dots,M_D)$, where each \(M_t\) is a partial matching on the active set $K.$ Let $J(\mathcal M):=\sum_{t=1}^D |M_t|$ be the number of two-qubit gate locations. Since each \(M_t\) is a partial
matching, 
\[J(\mathcal M)\le \frac{DQ}{2}\le Dn2^D.\]

Set $J_0:=Dn2^D.$ The unitary group \(U(4)\), equipped with the operator norm, has covering numbers
\[
    N(U(4),\|\cdot\|_{\mathrm{op}},\delta)
    \le
    \left(\frac{C}{\delta}\right)^C,
    \qquad 0<\delta\le1.
\]
Taking the product over all gate locations in the fixed architecture gives a
gate-net of cardinality at most
\[
    \left(\frac{C}{\delta}\right)^{C J(\mathcal M)}
    \le
    \left(\frac{C}{\delta}\right)^{C J_0}.
\]

We now relate the gate accuracy \(\delta\) to the residual metric. Let \(\theta_U\) and \(\theta_V\) be two circuit parameters with the same
architecture. Enumerate their corresponding gate locations in circuit order and write
\[
    U=A_J A_{J-1}\cdots A_1,
    \qquad
    V=B_JB_{J-1}\cdots B_1,
\]
where \(J=J(\mathcal M)\), and suppose that $\|A_j-B_j\|_{\mathrm{op}}\le\delta$ for $j=1,\dots,J$. The telescoping expansion gives
\[
    U-V
    =
    \sum_{j=1}^J
    A_J\cdots A_{j+1}
    (A_j-B_j)
    B_{j-1}\cdots B_1,
\]
and therefore $\|U-V\|_{\mathrm{op}} \le J\delta \le J_0\delta$. For any physical Pauli observable \(X=\sigma_i^b\otimes I_{\mathrm{anc}}\), we have
\[
    \|U^*XU-V^*XV\|_{\mathrm{op}}
    \le
    2\|U-V\|_{\mathrm{op}}
    \le
    2J_0\delta.
\]
Since all pulled-back Pauli observables are unitary, for
\(I=\{i_1,\dots,i_p\}\) and \(a=(a_1,\dots,a_p)\),
\[
\big\|
\prod_{s=1}^p O_{i_s}^{a_s}(U)
-
\prod_{s=1}^p O_{i_s}^{a_s}(V)
\big\|_{\mathrm{op}}
\le
\sum_{s=1}^p
\|O_{i_s}^{a_s}(U)-O_{i_s}^{a_s}(V)\|_{\mathrm{op}}
\le
2pJ_0\delta.
\]
Similarly, for every physical site \(i\in[n]\) and every \(b\in\{1,2,3\}\),
\[
\left|
\langle 0^{n+r}|O_i^b(U)|0^{n+r}\rangle
-
\langle 0^{n+r}|O_i^b(V)|0^{n+r}\rangle
\right|
\le
\|O_i^b(U)-O_i^b(V)\|_{\mathrm{op}}
\le
2J_0\delta .
\]
Since these expectations have absolute value at most \(1\), it follows that
\[
\big|
\prod_{s=1}^p
\langle 0^{n+r}|O_{i_s}^{a_s}(U)|0^{n+r}\rangle
-
\prod_{s=1}^p
\langle 0^{n+r}|O_{i_s}^{a_s}(V)|0^{n+r}\rangle
\big|
\le
2pJ_0\delta .
\]
Hence
\[
    |\chi_{I,a}(\theta_U)-\chi_{I,a}(\theta_V)|
    \le C_pJ_0\delta,
\]
and so
\[
\mathsf d_R(\theta_U,\theta_V)^2
=
\binom np^{-1}
\sum_{I\in\mcI_p^n}
\sum_{a\in\{1,2,3\}^p}
|\chi_{I,a}(\theta_U)-\chi_{I,a}(\theta_V)|^2
\le C_pJ_0^2\delta^2.
\]
Thus $\mathsf{d}_{R}(\theta_U,\theta_V)\le C_pJ_0\delta$. If \(0<\eta\le C_p\), choose $\delta=\frac{\eta}{C_pJ_0}.$ Then the product gate-net is an \(\eta\)-net in \(\mathsf{d}_{R}\) for a
fixed architecture, and its logarithmic size is at most $C_pJ_0\log (C_pJ_0/\eta)$.

If \(\eta>C_p\), then one representative per fixed architecture is enough
after increasing \(C_p\), because \(|\chi_{I,a}|\le2\) implies that the
\(\mathsf{d}_{R}\)-diameter is bounded by a constant depending only on $p$.
Thus, for all \(\eta>0\), the gate contribution for a fixed architecture is
bounded by
\[
    C_pJ_0\log_+\left(\frac{C_pJ_0}{\eta}\right)
    =
    C_pDn2^D
    \log_+\left(\frac{C_pDn2^D}{\eta}\right).
\]

It remains to combine the architecture count with the fixed-architecture gate nets. By the pruning and relabeling construction above, every circuit parameter
\(\theta\in\mathsf{T}_{n,r}^{(D)}\) is \(\mathsf{d}_{R}\)-equivalent to a
pruned, relabeled circuit parameter whose nonidentity gates act only on the
canonical active set \(K\), with \(|K|\le n2^D\). Thus it suffices to cover the collection of such pruned, relabeled circuits.

We showed that the number of possible pruned architectures on \(K\) has logarithm at most $C Dn2^D(\log n+D)$. For each fixed pruned architecture, the preceding gate-discretization argument
gives an \(\eta\)-net in \(\mathsf{d}_{R}\) of logarithmic size at most $C_pDn2^D
    \log_+ (C_pDn2^D/\eta)$. Taking the union of these fixed-architecture nets over all pruned
architectures gives a net for the full parameter space \(\mathsf{T}_{n,r}^{(D)}\).
Equivalently, the covering number is bounded by the number of pruned
architectures times the largest fixed-architecture covering number. Therefore,
after absorbing constants into \(C_p\),
\[
\log N(\mathsf{T}_{n,r}^{(D)},\mathsf{d}_{R},\eta)
\le
C_pDn2^D(\log n+D)
+
C_pDn2^D
\log_+\left(\frac{C_pDn2^D}{\eta}\right).
\]
The right-hand side is independent of $r$, as claimed.

\end{proof}

\end{document}
